\noindent\textbf{What this section should be:}
Explain how you summarize and compare results: what variability is measured, how error bars are computed, which statistical tests are used, and what assumptions they require. Prefer distributional comparisons when appropriate and justify meta-evaluation criteria.
\begin{itemize}
  \item \textit{[ACM]} ``enumerates and validates assumptions of statistical tests used (if any)''
  \item \textit{[ACM]} ``demonstrates appropriate statistical power (for quantitative work) or saturation (for qualitative work)''
  \item \textit{[ACM]} ``compares distributions (rather than means) of results using appropriate statistics''
  \item \textit{[ACM]} ``compares solutions using an appropriate meta-evaluation criteria; justifies the chosen criteria''
  \item \textit{[NeurIPS]} ``The factors of variability that the error bars are capturing should be clearly stated (for example, train/test split, initialization, random drawing of some parameter, or overall run with given experimental conditions).''
  \item \textit{[NeurIPS]} ``The method for calculating the error bars should be explained (closed form formula, call to a library function, bootstrap, etc.)''
  \item \textit{[NeurIPS]} ``The assumptions made should be given (e.g., Normally distributed errors).''
  \item \textit{[NeurIPS]} ``It should be clear whether the error bar is the standard deviation or the standard error of the mean.''
  \item \textit{[NeurIPS]} ``For asymmetric distributions, the authors should be careful not to show in tables or figures symmetric error bars that would yield results that are out of range (e.g. negative error rates).''
\end{itemize}